\chapter{Code and Reproducibility}
\label{app:code}

Every figure of the book is drawn by a Python script in the folder \textit{code} of the \LaTeX\ source. A figure is \emph{computed} when its script also produces its data, by a simulation or by evaluating a formula, \emph{illustrated} when its script draws a schematic, and \emph{replotted} when its script draws results that earlier experiments recorded in files. A replotted figure is a computed figure whose computation is kept apart from the drawing; the experiments are in the folder \textit{code/anc}, described in Section~\ref{sec:code_partiv}.

Unless an entry says otherwise, a script needs Python 3 with NumPy, SciPy and Matplotlib and no other package. The scripts \textit{figures.py}, \textit{figures\_color.py} and \textit{mf\_nls.py} call \texttt{numpy.trapezoid}, which NumPy introduced in version 2.0, so they need NumPy 2.0 or later. We ran the scripts of this appendix with Python 3.13.11, NumPy 2.4.2, SciPy 1.17.0 and Matplotlib 3.10.8 on a laptop with an Apple M2 processor.

\section{Running the scripts}
\label{sec:code_running}

Run a script from the folder \textit{code} with \texttt{python3} followed by its name, for example \texttt{python3 figures.py}. Unless its entry mentions options, a script takes no arguments. It writes its figures into the folder \textit{figures} next to \textit{code}, replacing the files shipped with the book, and prints the numbers that its entry lists; to keep the shipped figures, run the script in a copy of the source folder. A new figure enters the book at the next build (Section~\ref{sec:code_rebuild}).

Every script that draws random numbers seeds its generator with fixed integers, given in its entry. In \textit{figures.py} and in \textit{figures\_color.py} one generator serves all the figures of the file, so the figure of an entry depends on the random numbers that the entries before it used. The times in the entries are wall clock times that we measured on the laptop while other jobs were running on it.

Text and lines in the PDF files are vector graphics. The paperback prints them in grayscale; the digital edition, the PDF of the book's DOI, has them in color. Eleven files embed rasters at 300 pixels per inch: \textit{henon\_heiles}, \textit{cat\_map}, \textit{logistic}, \textit{sphere\_maxwell}, \textit{sinkhorn}, \textit{ns2d}, \textit{ns2d\_posterior}, \textit{ns2d\_posterior\_data}, \textit{markov\_switching}, \textit{particle\_learning} and \textit{fk\_posterior}.

Two figures show chaotic orbits, the section of the H\'enon Heiles system in Figure~\ref{fig:hh} and the standard map in Figure~\ref{fig:standard}. Nearby orbits separate exponentially, so a difference in the last bit of one floating point operation changes the position of a late point of an orbit by a macroscopic amount. A rerun on another machine, or with other versions of NumPy, SciPy or the mathematical library of the compiler, can therefore move individual points of these two figures, even with the same seed. Their structure does not change: the invariant curves, the islands and the chaotic seas return, and so does every number that the text quotes. The figures in the book were produced by these scripts on the authors' machine. The script \textit{hh\_checks.py} checks this behavior for Figure~\ref{fig:hh}. It integrates every orbit a second time with the relative tolerance $10^{-10}$ in place of $10^{-9}$, and it computes the largest Lyapunov exponent along each orbit from the tangent flow, renormalized every $5$ units of time up to $t=2500$. The ten orbits that start at $y\ge0$ agree with the first integration to within $10^{-5}$ over their first hundred crossings of the plane and never differ from it by more than $0.01$, while the five that start at $y\le-0.05$ differ from it by more than $0.01$ within their first hundred crossings. The exponents are at most $0.00027$ for the first ten and between $0.0316$ and $0.0614$ for the other five, so these are the ten orbits on invariant curves and the five chaotic orbits of the caption of the figure. The script needs NumPy and SciPy, draws no random numbers, writes no file and took between 111 and 168 seconds in our three runs. A rerun of the script for Figure~\ref{fig:standard} gave a picture with the same curves, islands and sea, in which about a fifth of the pixels, counted at $100$ pixels per inch, differ from the shipped file.

A reader with an older copy will find that the present \textit{figures.py} embeds images at 300 pixels per inch, adds a color bar to the H\'enon Heiles section and the energy and enstrophy spectra to the Navier Stokes figure, sums the P\'olya kernel over thirty terms where the older copy summed twenty nine and locates the zeros of $H_t$ by a root solve, that the present \textit{figures\_color.py} adds color bars to the logistic map and the sphere, adds the limit curve to the Lanford window and computes the orbit of the doubling map as a shift of binary digits where the older copy drew independent uniform numbers, and that \textit{figures\_pf.py} and \textit{figures\_msw.py} now print the numbers that the text quotes.

\section{Figures, tables and scripts}
\label{sec:code_figures}

Table~\ref{tab:code_figures} lists the figures in the order of the book. Table~\ref{tab:code_tables} lists the tables whose numbers a script of the folder computes or reads.

\begin{table}[tp]\centering\footnotesize\setlength{\tabcolsep}{4pt}
\caption{The figures of the book, the files that hold them and the scripts that write them.}\label{tab:code_figures}
\begin{tabular}{llll}
\toprule
Figure & File & Script & Kind\\
\midrule
\ref{fig:hh} & \textit{henon\_heiles.pdf} & \textit{figures.py} & computed\\
\ref{fig:birkhoff} & \textit{birkhoff.pdf} & \textit{figures\_color.py} & computed\\
\ref{fig:cat} & \textit{cat\_map.pdf} & \textit{figures\_color.py} & computed\\
\ref{fig:logistic} & \textit{logistic.pdf} & \textit{figures\_color.py} & computed\\
\ref{fig:standard} & \textit{standard\_map.pdf} & \textit{figures\_color.py} & computed\\
\ref{fig:types} & \textit{types\_simplex.pdf} & \textit{figures\_color.py} & computed\\
\ref{fig:conditioning} & \textit{conditioning.pdf} & \textit{figures\_color.py} & computed\\
\ref{fig:spheremaxwell} & \textit{sphere\_maxwell.pdf} & \textit{figures\_color.py} & computed\\
\ref{fig:regret} & \textit{regret.pdf} & \textit{figures\_color.py} & computed\\
\ref{fig:exit} & \textit{exit\_time.pdf} & \textit{figures.py} & computed\\
\ref{fig:lanford} & \textit{lanford\_window.pdf} & \textit{figures\_color.py} & illustrated\\
\ref{fig:kac} & \textit{kac\_model.pdf} & \textit{figures.py} & computed\\
\ref{fig:kacring} & \textit{kac\_ring.pdf} & \textit{figures.py} & computed\\
\ref{fig:heat_inverse} & \textit{heat\_inverse.pdf} & \textit{figures.py} & computed\\
\ref{fig:fk_posterior} & \textit{fk\_posterior.pdf} & \textit{fk\_posterior.py} & computed\\
\ref{fig:ns2d} & \textit{ns2d.pdf} & \textit{figures.py} & computed\\
\ref{fig:ns2d_posterior_data} & \textit{ns2d\_posterior\_data.pdf} & \textit{ns2d\_posterior.py} & computed\\
\ref{fig:ns2d_posterior} & \textit{ns2d\_posterior.pdf} & \textit{ns2d\_posterior.py} & computed\\
\ref{fig:burgers} & \textit{burgers.pdf} & \textit{figures.py} & computed\\
\ref{m:fig:radius} & \textit{mf\_radius\_theory.pdf} & \textit{mf\_radius\_theory.py} & replotted\\
\ref{m:fig:heat} & \textit{mf\_heat.pdf} & \textit{mf\_heat.py} & computed\\
\ref{m:fig:heatgbc} & \textit{mf\_heat\_gbc.pdf} & \textit{mf\_heat\_gbc.py} & replotted\\
\ref{m:fig:tg} & \textit{mf\_tg.pdf} & \textit{mf\_tg.py} & computed\\
\ref{m:fig:cost} & \textit{mf\_cost.pdf} & \textit{mf\_cost.py} & replotted\\
\ref{m:fig:nls} & \textit{mf\_nls.pdf} & \textit{mf\_nls.py} & computed\\
\ref{fig:lee_yang} & \textit{lee\_yang.pdf} & \textit{figures.py} & computed\\
\ref{fig:sinkhorn} & \textit{sinkhorn.pdf} & \textit{figures.py} & computed\\
\ref{fig:pf} & \textit{pf\_errors.pdf} & \textit{figures\_pf.py} & computed\\
\ref{fig:pl} & \textit{particle\_learning.pdf} & \textit{figures\_pl.py} & computed\\
\ref{fig:msw} & \textit{markov\_switching.pdf} & \textit{figures\_msw.py} & computed\\
\ref{fig:kelly} & \textit{kelly.pdf} & \textit{figures.py} & computed\\
\ref{fig:debruijn_newman} & \textit{debruijn\_newman.pdf} & \textit{figures.py} & computed\\
\bottomrule
\end{tabular}
\end{table}

\begin{table}[t]\centering\footnotesize\setlength{\tabcolsep}{4pt}
\caption{The tables of the book whose numbers come from a script of the folder \textit{code}.}\label{tab:code_tables}
\begin{tabular}{lp{2.2in}p{1.6in}}
\toprule
Table & Script & Numbers\\
\midrule
\ref{tab:yor_checks} & \textit{yor\_catalogue.py} & computed by the script\\
\ref{tab:fk_posterior} & \textit{fk\_posterior.py} & computed by the script\\
\ref{tab:ns2d_posterior} & \textit{ns2d\_posterior.py} & computed by the script\\
\ref{m:tab:gain} & \textit{mf\_gain.py} & computed by the script\\
\ref{m:tab:gap} & \textit{mf\_gap.py} & computed by the script\\
\ref{m:tab:iota} & \textit{mf\_heat.py} & computed by the script\\
\ref{m:tab:diag} & \textit{exp\_paper\_figs.py} & \textit{results/paper\_figs.json}\\
\ref{m:tab:cost} & \textit{exp\_paper\_figs.py} & \textit{results/paper\_figs.json}\\
\ref{m:tab:cover} & \textit{exp\_cover2.py}, \textit{agg\_cover2.py}; the columns $z$ and $p$ by \textit{mf\_cost.py} & \textit{results/cover2\_agg.json} and \textit{results/cover2\_rep*.json}\\
\ref{m:tab:ns} & \textit{variants.py}, \textit{hetero.py}, \textit{centring.py}, \textit{correct\_centre.py}, \textit{mcmc\_error.py} of the folder \textit{ns}; the row of the rescaled spread by \textit{table\_ablation.py} of the folder \textit{ns\_checks} & \textit{results/ns}; that row computed by the script\\
\bottomrule
\end{tabular}
\end{table}

\section{The figure scripts}
\label{sec:code_scripts}

The script \textit{figures.py} draws eleven figures and takes between 40 and 90 seconds. It uses one generator, seeded with $1900$.

\begin{description}
\item[H\'enon Heiles.] The initial $p_x>0$ of each orbit is fixed by the energy, and each orbit is integrated to time $2500$ by the default method of SciPy's \texttt{solve\_ivp}, an explicit Runge Kutta method with step size control, with relative tolerance $10^{-9}$ and absolute tolerance $10^{-11}$. Event detection finds the crossings of the plane $x=0$ with $p_x>0$, which give between $362$ and $412$ points per orbit in the plane $(y,p_y)$ (Figure~\ref{fig:hh}). The script \textit{hh\_checks.py} integrates the same orbits for the checks of Section~\ref{sec:code_running}.
\item[Exit time.] Six thousand Brownian paths started at zero are simulated by the Euler scheme with step $2\times10^{-4}$ until they leave $(-1,1)$. Sixty thousand draws of the series $\sum_n8E_n/((2n-1)^2\pi^2)$ of Theorem~\ref{thm:exit_series}, with independent standard exponential $E_n$ and two hundred terms, give the second histogram, and the density of the theorem, summed over two hundred terms, is the dashed curve (Figure~\ref{fig:exit}).
\item[Kac's model.] Twenty thousand velocities start from an equal mixture of two Gaussian peaks at $\pm1$ with standard deviation $0.1$, rescaled to unit second moment, and are updated by rotating random pairs through random angles, fifteen units of time at $N/2$ rotations per unit. The kurtosis is recorded every tenth of a unit (Figure~\ref{fig:kac}).
\item[Kac's ring.] A ring of $500$ sites with $61$ marked edges, chosen at random, is run for $2N+50$ steps from the state with all balls white. The figure shows (white minus black)$/N$ against Boltzmann's prediction $(1-2\mu)^t$ with $\mu=61/500$ (Figure~\ref{fig:kacring}).
\item[Heat equation.] Forty sine modes on $[0,\pi]$ with prior standard deviations $k^{-3/2}$, observation time $T=0.05$ and noise level $0.002$. The true coefficients are drawn from the prior, the posterior of each mode is computed in closed form, and two hundred posterior draws give the band from the $5$th to the $95$th percentile (Figure~\ref{fig:heat_inverse}).
\item[Lee Yang zeros.] The partition function of the Ising model on a $4\times4$ torus at $\beta J=0.6$, with $J=1$, is computed exactly as a polynomial of degree sixteen in the fugacity by enumerating all $2^{16}$ configurations, and its roots are found by the companion matrix method (\texttt{numpy.roots}). The script prints the largest deviation of a root modulus from one, $1.2\times10^{-15}$, which the text and the caption of Figure~\ref{fig:lee_yang} quote. The number is a rounding error, so its digits can change with the machine and the linear algebra library.
\item[de Bruijn Newman.] The P\'olya kernel $\Phi(u)$ is summed over thirty terms on $20{,}001$ equally spaced points of $[0,1.4]$, and $H_t(x)=\int_0^{1.4}e^{tu^2}\Phi(u)\cos(xu)\dd u$ is computed by the trapezoidal rule at $900$ equally spaced points of $[18,56]$ for $t=-4$, $0$ and $4$. The script finds the first three zeros of each $H_t$ as sign changes on that grid and refines them by SciPy's \texttt{brentq} applied to the same quadrature. It prints them to two decimals, so that $45.20$ appears as $45.2$: $30.62$, $45.2$, $52.29$ for $t=-4$; $28.27$, $42.04$, $50.02$ for $t=0$; $25.99$, $39.14$, $47.64$ for $t=4$ (Figure~\ref{fig:debruijn_newman}). The caption says that the zeros for $t=0$ are twice the ordinates of the zeros of the zeta function. We confirmed this against the first three ordinates computed with the package mpmath, which no script of the book uses: they are $14.13$, $21.02$ and $25.01$ to two decimals, and half of the zeros of $H_0$ agree with them to within $5\times10^{-6}$.
\item[Entropic transport.] Sinkhorn iterations, three thousand for each value of $\varepsilon$, between a law with two Gaussian shaped peaks and a law with one, both on $160$ points of $[-4,4]$, with the kernel $e^{-|x-y|^2/\varepsilon}$ and $\varepsilon=2$, $0.3$ and $0.03$, the values in the panel titles (Figure~\ref{fig:sinkhorn}).
\item[Kelly.] The growth rate $p\log_2(1+f)+(1-p)\log_2(1-f)$ in bits per bet for $p=0.6$ and $0\le f\le0.6$, with the optimal fraction $f^*=2p-1$, the maximum $1-H(p)$ and the point $2f^*$ marked (Figure~\ref{fig:kelly}). The script prints the maximum, $0.0290$ bits, the zero of the growth rate found by \texttt{brentq}, $f=0.38939$, which the caption quotes as $0.389$, and the growth rate at $2f^*=0.4$, $-0.00353$ bits.
\item[Burgers.] The viscous Burgers equation with initial condition $u_0=\sin x$ is solved by the Hopf Cole transform for $\nu=0.2$, $0.05$ and $0.005$ at $t=1.5$. The integral of the transform is evaluated by quadrature on $6000$ points of $[-\pi,3\pi]$, at $400$ equally spaced points of $[0,2\pi]$ and at $1201$ more points within $0.15$ of $x=\pi$, where the shock forms (Figure~\ref{fig:burgers}).
\item[Two dimensional Navier Stokes.] A pseudospectral solution of the vorticity equation on the torus $[0,2\pi]^2$ with $96\times96$ grid points, the two thirds rule for dealiasing, viscosity $2\times10^{-3}$ and fourth order Runge Kutta steps of $2\times10^{-3}$ up to $t=5$. The initial vorticity is a sum of ten Gaussian vortices $s\,e^{-d^2/0.12}$, where $d$ is the distance on the torus to a uniformly random center and the strength $s$ has a random sign and an absolute value uniform between $3$ and $6$ (Figure~\ref{fig:ns2d}). The script prints the largest absolute value of the initial vorticity, $5.87$, and for each of the three times the total energy and enstrophy and summaries of the two spectra. Section~\ref{sec:simns2d} quotes the changes of the total energy and enstrophy, the energy in the first shell, the share of the energy at wavenumbers up to three, the share of the enstrophy at wavenumbers from sixteen and the mean wavenumber of the enstrophy.
\end{description}

The script \textit{figures\_color.py} draws nine figures and takes between 7 and 11 seconds. It uses one generator, seeded with $2026$, and a second generator, seeded with $7$, for the last digits of the orbit of the doubling map.

\begin{description}
\item[Cat map.] Arnold's cat map acts on a color image on a $101\times101$ grid. The script iterates until the image returns, which takes $25$ steps and is printed, and draws the images after $0$, $1$, $2$, $5$, $12$ and $25$ steps (Figure~\ref{fig:cat}).
\item[Birkhoff averages.] Averages of $\cos2\pi x$ over the first $100{,}000$ points of an orbit of the rotation by the golden mean started at $0.1$, and of an orbit of the doubling map, on logarithmic axes against the rates $1/n$ and $1/\sqrt n$. The binary digits of the start of the doubling orbit are fair coin flips: the first $100{,}000$ come from the shared generator and the last $52$ from the second one, which leaves the random numbers of the later figures unchanged (Figure~\ref{fig:birkhoff}).
\item[Standard map.] Ninety orbits of the Chirikov standard map with $K=0.971635$, each of $900$ points from a random start on the torus $[0,2\pi)^2$ and each drawn in its own color (Figure~\ref{fig:standard}). The chaotic orbits are sensitive to rounding, as described in Section~\ref{sec:code_running}.
\item[Logistic map.] For $1600$ values of $r$ from $2.8$ to $4$ the script iterates $x\mapsto rx(1-x)$ from $x=0.3$, discards $600$ iterates, averages $\log|r(1-2x)|$ over the next $400$ to estimate the Lyapunov exponent $\chi(r)$, and plots the last $150$ of these iterates colored by $\chi$ clipped to $[-0.8,0.7]$ (Figure~\ref{fig:logistic}).
\item[Types on the simplex.] The contours are $\KL(Q\,\|\,P)$ for $P=(0.2,0.3,0.5)$ on a $300\times300$ grid, the line is $\{Q:\E_Qf=0.6\}$ for $f(x)=x$, and the I projection $P^*$ is the exponential tilt of $P$ whose multiplier is found with \texttt{brentq}. The script prints $P^*=(0.5587,0.2825,0.1587)$ and the multiplier $-1.0874$ (Figure~\ref{fig:types}).
\item[Moment conditioning.] The density of $X_1$ given that the mean of $n$ independent exponential variables equals $a$ is a rescaled Beta$(1,n-1)$ density, drawn for $a=2$ and $n=3,5,10,50$ against the exponential density of mean $a$. Its total variation distance to that density is computed by the trapezoidal rule for $29$ values of $n$ from $3$ to $1995$ and plotted on logarithmic axes. The script prints $n$ times this distance at the largest $n$, $0.3101$, close to the constant $c\approx0.31$ of the text (Figure~\ref{fig:conditioning}).
\item[Uniform points on a sphere.] For $n=3$, $6$ and $60$ the script draws $200{,}000$ points uniformly on the sphere of radius $\sqrt n$ in $\R^n$, as normalized Gaussian vectors, and plots the joint histogram of their first two coordinates on $120\times120$ bins (Figure~\ref{fig:spheremaxwell}).
\item[Lanford window.] The left panel plots the partial sums of the geometric majorant $\sum_k(t/T^*)^k$ for $k\le2$, $5$, $10$ and $40$ against its limit $1/(1-t/T^*)$, and the right panel the bound $r^k$ on the $k$th term for $t=rT^*$ with $r=0.5$, $0.9$ and $1.1$. The script evaluates these formulas and solves no equation (Figure~\ref{fig:lanford}).
\item[Regret of Bayes mixtures.] The worst case over $k$ of the logarithm of the ratio of the maximized Bernoulli likelihood to the Beta$(a,a)$ mixture probability, in bits, for $n=2,\dots,2000$ and $a=1$ (Laplace) and $a=\frac12$ (Jeffreys), computed with log gamma functions and plotted against $\frac12\log_2(\pi n/2)$ and $\log_2(n+1)$ (Figure~\ref{fig:regret}).
\end{description}

The scripts \textit{figures\_pf.py}, \textit{figures\_msw.py} and \textit{figures\_pl.py} draw one figure each for Chapter~\ref{ch:filters}, with the seeds $1993$, $7$ and $2010$, and take about 3, 1 and 2 seconds.

\begin{description}
\item[Particle filter error.] A local level model with $q=0.1$, $r=1$ and $T=200$ observations is simulated and filtered exactly by the Kalman filter and by the bootstrap particle filter (Figure~\ref{fig:pf}). The script prints the least squares slopes of the variance over sixty runs against $t$, $0.0125$ for $N=100$ and $0.00117$ for $N=1000$; the effective sample size without resampling at $t=10$, $20$, $30$ and $200$, $181.5$, $25.8$, $3.7$ and $1.0$; its minimum with resampling, $350.3$; and the numbers of distinct ancestors at $t=1$, $10$, $50$ and $200$, $570$, $105$, $20$ and $6$. The text quotes the ratio of the slopes as a factor of about ten, the effective sample sizes and the numbers of ancestors.
\item[Markov switching filter.] The two regime model of the text is simulated for $600$ steps and filtered by $4000$ particles (Figure~\ref{fig:msw}). The script prints the accuracy of the regime classification, $0.933$; the times of the two switches, $4.95$ and $14.35$; the misclassified steps, $40$ of $599$; and the delay after which the filter agrees with the new regime, $8$ steps ($0.4$ time units) and $13$ steps ($0.65$). The text quotes the accuracy and the two delays.
\item[Particle learning.] The local level model of Example~\ref{ex:pl_local_level}, with the settings of the text, is filtered by particle learning and by a naive bootstrap filter (Figure~\ref{fig:pl}). The script prints the final $5$th, $50$th and $95$th percentiles of $q$, $0.2206$, $0.2888$ and $0.3468$, and of $r$, $0.8804$, $1.0149$ and $1.1732$, and the number of distinct values of $q$ left in the naive filter, $3$. The text quotes the two medians and the count.
\end{description}

The script \textit{fk\_posterior.py} draws one figure and one table for Chapter~\ref{ch:nsinference}.

\begin{description}
\item[Fisher Kolmogorov posterior.] The script evaluates the posterior of $(\log D,\log r)$ exactly on a grid and with no Markov chain, for three nested observation windows $[0,T]$ (Figure~\ref{fig:fk_posterior} and Table~\ref{tab:fk_posterior}). The script prints the largest difference of the two solvers of the chapter on the data grid, below $6\times10^{-6}$ in the checks it makes, the speed of the front, the posterior of each window, a linearized analysis at the truth, a calibration study with $1000$ other noise draws, and a study of a small conditional generator written in PyTorch. The seeds are $1937$ for the data, $2026$ for the calibration draws, $11$ and $12$ for the simulated pairs of the generator study and $3$, $4$ and $5$ for its training. The script uses at most three threads and a full run takes about ten minutes; we measured $8.9$ and $12.7$ minutes. An option that skips the generator study takes about three minutes ($2.5$ minutes in our run) and does not need PyTorch, which is used only there (we used version 2.11.0). The environment variable \texttt{FK\_OUT} names another folder for the figure, and \texttt{FK\_CACHE} names a file that stores the results, from which another option redraws the figure in a second. Our full run printed the same numbers as the run that produced the shipped figure and wrote a figure file that differs from the shipped one only in the creation date. The text of Chapter~\ref{ch:nsinference} quotes numbers that the script prints, and we compared the entries of Table~\ref{tab:fk_posterior}, the speed of the front, the calibration percentages and the numbers of the generator study with the printed lines.
\end{description}

The script \textit{ns2d\_posterior.py} draws two figures and computes one table for Chapter~\ref{ch:nsinference}. It needs JAX in addition to NumPy, SciPy and Matplotlib; we used JAX 0.9.1, on the processor and in double precision.

\begin{description}
\item[Posterior of the initial vorticity.] The script finds the maximum of the posterior by a Levenberg Marquardt iteration, samples the posterior by Hamiltonian Monte Carlo with four chains, with the gradient computed as the reverse mode derivative of the solver, and repeats the sampling with shorter chains for the noise levels $0.125$ and $0.5$ (Figures~\ref{fig:ns2d_posterior_data} and~\ref{fig:ns2d_posterior} and Table~\ref{tab:ns2d_posterior}). It prints, on lines that start with RESULT, the numbers that Section~\ref{sec:simns2d} quotes about the posterior: the convergence diagnostics, the ratios of posterior to prior standard deviation by wavenumber shell, the correlations of the coefficients, the entries of the table and the results at the other noise levels. Its seeds are $2026$ for the noise, $1$ for the start of the first optimization and $101$, $102$ and $103$ for three more, $7$ for the starting points of the chains, $11$ for the chains, $3$ and $6$ for the prior draws behind the prior standard deviation of the energy and the position errors under the prior, $5$ for the draws behind the importance weights and $99$ for the checks. A full run uses three processes and took 19 minutes. It keeps the chains and the intermediate results in a folder, by default in the temporary folder of the system and otherwise in the one given after the option \texttt{-{}-cache}. The option \texttt{-{}-plot-only} reads that folder, prints the numbers and redraws the figures in about ten seconds, and \texttt{-{}-quick} runs short chains for testing, which took seven minutes. The option \texttt{-{}-checks}, which took six minutes, compares the gradient with central differences and runs the sampler on the problem with the nonlinear term dropped, whose posterior is Gaussian and known exactly. It prints the largest $\widehat R$ of that run, $1.009$, and the ratios of the sampled to the exact standard deviations, between $0.946$ and $1.057$, which the section quotes as agreement to within $6$ percent. Our full run printed the numbers that the section quotes, with position errors under the prior between $0.717$ and $0.729$ for the five vortices, and the figures it wrote are identical to the shipped files except for the creation date.
\end{description}

The script \textit{yor\_catalogue.py} draws no figure and computes Table~\ref{tab:yor_checks} for the section on Yor's identities as algorithms in Chapter~\ref{ch:brownian}.

\begin{description}
\item[Yor identities.] The script needs NumPy and SciPy only. For each exact sampler of Table~\ref{tab:yor_catalogue} that the section tests, it prints the Kolmogorov Smirnov distance of $10^5$ draws to the exact distribution function, and the same distance for the Euler schemes, which are the entries of Table~\ref{tab:yor_checks}. It also prints the calibration of these distances over $40$ repetitions with other seeds, the checks of the Kent series, of the identity of Matsumoto and Yor, of the moments of Biane, Pitman and Yor and of the horseshoe law, and the times of the exact samplers and of the Euler schemes that the section quotes. Its seeds are $20261007$ plus a fixed offset for each experiment. It writes its results into the file \textit{yor\_results.json} of a folder \textit{yor\_out} in the folder from which it is run, or of the folder given after the option \texttt{-{}-out}, and the option \texttt{-{}-fast} runs a test with ten times fewer samples, which took 10 seconds. A full run took 74 seconds on a lightly loaded machine and 196 seconds on a busy one. The two runs printed the same numbers except the times and their ratios, and these numbers agree with the section.
\end{description}

\section{Experiments on the solution manifold}
\label{sec:code_partiv}

The six scripts \textit{mf\_radius\_theory.py}, \textit{mf\_heat.py}, \textit{mf\_heat\_gbc.py}, \textit{mf\_tg.py}, \textit{mf\_cost.py} and \textit{mf\_nls.py} draw the figures of Part~\ref{part:manifold}. They use no random numbers, each runs in under a second, and each prints the values that the text quotes. Five of them were rebuilt from four figure scripts of the folder \textit{code/anc}, described below, with the formulas and numbers of those scripts: \textit{mf\_heat.py} and \textit{mf\_tg.py} from \textit{fig\_heat\_stokes.py}, \textit{mf\_heat\_gbc.py} from \textit{fig\_heat\_gbc.py}, \textit{mf\_cost.py} from \textit{fig\_cost.py} and \textit{mf\_nls.py} from \textit{fig\_nls\_bogoliubov.py}. The script that drew the radius figure of the paper is not in the folder, so \textit{mf\_radius\_theory.py} rebuilds that figure from the recorded measurements. Two further scripts, \textit{mf\_gain.py} and \textit{mf\_gap.py}, were written for the book from the settings of Tables~\ref{m:tab:gain} and~\ref{m:tab:gap}. They print these tables and draw no figure, need only NumPy, use no random numbers and run in under a second.

\begin{description}
\item[Decomposition of the radius error.] \textit{mf\_radius\_theory.py} fits the decomposition of the error against the localization radius to seven measured errors for each of two heads (Figure~\ref{m:fig:radius}). The measurements are those of \textit{anc/results/ns/radius.json}, copied into the script, which therefore runs without the folder \textit{anc}. The script prints the fitted constants, the root mean square residuals $0.0646$ for the identity head and $0.0440$ for the affine head, and the optimal radius of the affine head, $c^*=2.44$. The measurements come from \textit{radius.py} of the Navier Stokes experiment below, which we reran.
\item[Information per mode.] \textit{mf\_heat.py} evaluates the information $\iota_k(T_0)$ per mode of the heat equation on the torus with $\nu=0.01$ and $T=1$ for $T_0=0$, $0.05$, $0.2$ and $0.5$, and the posterior standard deviation of the state against $t$ for $N=10^4$ and the truncations $K_{\max}=16$, $32$, $64$ and $128$ of the wave vectors in the plane (Figure~\ref{m:fig:heat}). It prints the mode at which each $\iota_k$ falls below $10^{-24}$, the plateau values at $t=10^{-5}$ and their ratios, the time from which each of the first three truncations agrees with the largest to within one percent, the common value at $t=1$, and the rows of Table~\ref{m:tab:iota}, which agree with the printout \textit{run\_log\_heat.txt} recorded in the folder \textit{anc}.
\item[Finite $N$ check.] \textit{mf\_heat\_gbc.py} draws the posterior of the heat equation on $\T^1$ with $N=2000$ and twenty modes for the prior scale $\rho=1$ and for the calibrated $\rho=124.03$, with the closed form prediction of the Bernstein von Mises standard deviations, from the recorded file \textit{anc/results/paper\_figs.json} (Figure~\ref{m:fig:heatgbc}). For the two settings it prints the number $k_{\rm eff}$ of modes on which the data dominate the prior, the share of the modes whose truth lies within two posterior standard deviations of the mean, the largest standardized error and the largest relative deviation of the exact standard deviations from the prediction, which are entries of Table~\ref{m:tab:diag}.
\item[Taylor Green coupling.] \textit{mf\_tg.py} evaluates the ratio of the advective coupling to the viscous decay of mode $k$ for the Navier Stokes equation linearized about the Taylor Green vortex, with $\nu=0.01$ and $A=0.1$, $1$ and $5$, and prints the wavenumber $|k|=A/(2\pi\nu)$ at which each ratio equals one, $1.59$, $15.92$ and $79.58$ (Figure~\ref{m:fig:tg}).
\item[Error against cost.] \textit{mf\_cost.py} draws the recorded error of GBC and of pCN against the number of solves, and the recorded coverage over replicates, from \textit{anc/results/study2\_agg.json} and \textit{anc/results/cover2\_agg.json} (Figure~\ref{m:fig:cost}). It prints the fitted slopes, $-0.521$ for GBC and $-0.507$ for pCN, the band of accuracy that both reach, $0.0811$ down to $0.0527$, and the ratio of the numbers of solves they need, $1154$ to $1181$, and the same band and ratio for the relative Frobenius error of the covariance, $0.1397$ down to $0.0628$ and $366$ to $510$. It also reads the twenty recorded replicate files \textit{anc/results/cover2\_rep1.json} to \textit{cover2\_rep20.json}, checks that they reproduce the coverages and the discordant counts of \textit{cover2\_agg.json}, and recomputes the paired test of Table~\ref{m:tab:cover} with the replicate as its unit, since \textit{agg\_cover2.py} treats the $2400$ (replicate, mode) pairs of an arm as independent. For the three truth arms it prints the discordant counts, the clustered statistics $z$ and the exact $p$ values of that table: $5$ and $3$ with $z=0.71$ and $p=0.73$, $2$ and $6$ with $z=-1.63$ and $p=0.22$, and $1$ and $1$ with $z=0$ and $p=1$. The coverage columns of that table are the entries of \textit{cover2\_agg.json}, which we compared with the table.
\item[Bogoliubov blocks.] \textit{mf\_nls.py} computes by the trapezoidal rule the information blocks $\Theta_p(T)$ of the cubic Schr\"odinger equation linearized about a plane wave, for the free equation and for $\kappa a^2=\pm20$, and draws the eigenvalues of $\Theta_p(1)$ against $|k|$ and the largest eigenvalue against $T$ for three modes of the unstable band (Figure~\ref{m:fig:nls}). It prints the free block at $T=1$, the eigenvalues at $|k|=1$, $2$, $4$, $8$ and $16$ for the defocusing and focusing cases, which agree with the printout \textit{run\_log\_nls.txt} recorded in the folder \textit{anc}, and the condition $\lambda_1=39.4784<2|\kappa|a^2=40$ that puts the lowest mode in the unstable band.
\item[Manifold gain.] \textit{mf\_gain.py} computes Table~\ref{m:tab:gain} with the settings of its caption: the spatial modes $0<\max(|k_1|,|k_2|)\le12$ in both priors, the temporal frequencies $|m|\le40$ in the unconstrained one, $\rho=1$ in the constrained prior, and in the unconstrained prior the scale $\rho=0.6928$ under which both priors give the signal of an observation the same prior variance, the average over $0\le t\le T$ of the pointwise prior variance of $u(t,x)$. It prints the two scales, that common prior standard deviation, $1.748$, the pointwise prior standard deviations, from $2.158$ at $t=0$ to $1.387$ at $t=1$ for the constrained prior and $1.748$ at every $t$ for the unconstrained one, the numbers of effective parameters, $20$ and $116$, and for $t=0$, $0.05$, $0.2$, $0.5$ and $1$ the posterior standard deviations, from $0.226915$ down to $0.051404$ for the constrained prior and $0.313350$ for the unconstrained one, and their ratios, $1.3809$, $1.6164$, $2.2902$, $3.6869$ and $6.0959$. With \texttt{MATCHED = False} it takes $\rho=1$ in both priors, the setting of the table of our paper~\cite{polsonsokolov2026manifold}, and prints the pointwise prior standard deviations $2.158$ and $2.523$, $20$ and $172$ effective parameters, the unconstrained standard deviation $0.377444$ and the ratios $1.6634$, $1.9470$, $2.7587$, $4.4411$ and $7.3427$; these agree to a relative difference of at most $2.3\times10^{-16}$ with the file \textit{manifold\_gain.json} that the paper recorded, which is not part of the book source.
\item[Sufficiency gap.] \textit{mf\_gap.py} computes Table~\ref{m:tab:gap} from the population information $\diag(N\iota_j)$ of the $288$ modes with $0<|k|_\infty\le8$, for the values $4$, $8$, $16$, $20$, $32$, $48$ and $288$ of $J'$ in its rows. It prints $D=288$, $J_{\rm eff}=20$ and the four columns of the table, and also the lower bound of Theorem~\ref{m:thm:gap}, which differs from the Rao Blackwell column by at most $10.7$ percent, and the ratio of the upper bound to that column, at most $3.53$, which the caption and the text quote as at most $11$ percent and within a factor of four.
\end{description}

The folder \textit{code/anc} is a copy of the ancillary files of our paper~\cite{polsonsokolov2026manifold}, where the experiments of Chapter~\ref{ch:examples} were first reported. It holds the Python scripts that ran them, the results that they recorded as JSON files in the folders \textit{results} and \textit{results/ns}, two batch scripts for the SLURM scheduler of a cluster, on which the replicate studies were run, and a file \textit{README.txt} that lists what produces what. The file \textit{README.txt} and the script \textit{ns/run\_all.sh} predate the book; where they differ from the folder, by naming files that it does not hold or other output folders, the entries below describe the folder as shipped. The seeds are fixed in the scripts. The replicate studies take hours on one core, so the figures and tables of the book read the recorded results and do not rerun them. The batch scripts change into a scratch directory of the authors' cluster account, which a reader has to replace.

\begin{description}
\item[Finite $N$ check and cost.] \textit{exp\_paper\_figs.py}, with the seeds $1$ and $7$, writes \textit{results/paper\_figs.json}, from which come Tables~\ref{m:tab:diag} and~\ref{m:tab:cost} and the data of Figure~\ref{m:fig:heatgbc}. It takes about 160 seconds. A rerun reproduces every entry of the file except the wall clock times and the effective sample size per second that is computed from them. The wall clock times are the column of Table~\ref{m:tab:cost} headed Time, which therefore depends on the machine: the table has the recorded $2.91$ seconds for GBC and $86.5$ seconds for pCN, and our rerun on a busy machine took $5.50$ and $151.8$.
\item[Error against cost.] \textit{exp\_study2.py} runs one replicate of the comparison of GBC and pCN for $N=50000$ observations and eight modes. Its seeds are $90000$ plus the index of the replicate for the data, $500000$ plus $1000$ times the index plus the budget (and an offset for each variant) for the generators, and $700000$ plus $1000$ times the index plus the number of iterations for the pCN chains. The forty replicates were run as an array of SLURM jobs (\textit{study2.sbatch}), and \textit{agg\_study2.py} combines the forty recorded files into \textit{results/study2\_agg.json}, which \textit{mf\_cost.py} reads. We reran the forty replicates from a copy of the folder as an array of jobs on a cluster with Intel Xeon processors, four threads to a job. A replicate took between 12 and 20 minutes, almost all of it in the pCN chains, so the forty take about ten hours in sequence. Every rerun file reproduces the recorded one to a relative difference of at most $6.7\times10^{-11}$ in every entry except the wall clock times, and so does the aggregate that \textit{agg\_study2.py} makes from them; from the forty recorded files it reproduces \textit{study2\_agg.json} exactly.
\item[Coverage over replicates.] \textit{exp\_cover2.py} runs one task of fifteen replicates for each of three truth arms, and the twenty tasks of \textit{cover2.sbatch} give the $300$ replicates of Table~\ref{m:tab:cover} and of the lower panel of Figure~\ref{m:fig:cost}. The seed of a replicate is $300000$ plus $1000$ times the index of the task plus the index of the replicate within the task, plus an offset of $0$, $400000$ or $800000$ for the arm, and the generator is seeded with that number plus $500000$; the fixed typical truth is drawn with the seed $20260826$. \textit{agg\_cover2.py} combines the twenty recorded files into \textit{results/cover2\_agg.json}, which \textit{mf\_cost.py} reads. We reran the twenty tasks in the same way, and a task took about two and a half minutes. Every rerun file, and the aggregate made from them, equals the recorded one exactly. From the rerun files and aggregates of the two studies \textit{mf\_cost.py} prints the same numbers and draws the same figure as from the recorded ones.
\item[Navier Stokes.] The folder \textit{ns} holds the pseudospectral solver \textit{ns\_solver.py} and the scripts of the experiment of Section~\ref{m:sec:expns}, which are run from the folder \textit{ns} in the order of \textit{run\_all.sh}; according to the \textit{README.txt} of the folder \textit{anc} the whole sequence takes about forty minutes on one core, and according to the header of \textit{run\_all.sh} forty five. The scripts that write results write one JSON file each, in the folder \textit{ns} itself (\textit{radius.json} for \textit{radius.py}, \textit{variants.json} for \textit{variants.py} and so on), except \textit{advection\_ratio.py}, \textit{mcmc\_error.py} and \textit{crossover.py}, which write \textit{advection.json}, \textit{mcmc\_error.json} and \textit{crossover2.json} or \textit{crossover3.json} into \textit{results/ns}, where the recorded copies of all the files are kept, so that a run of these three replaces the recorded copies. The script \textit{full.py} runs the reference Metropolis chain and the training pools and writes them to \textit{cache.npz}, which most of the others read; it took about eight minutes in our run. The script \textit{variants.py} computes the upper panel of Table~\ref{m:tab:ns}, \textit{hetero.py} (about five minutes) writes the residuals that \textit{centring.py} and \textit{correct\_centre.py} use for the rows of the lower panel, and \textit{mcmc\_error.py} computes the last row, the Monte Carlo error of the reference, $0.063$. The scripts \textit{radius.py} and \textit{radius\_iw.py} sweep the localization radius without and with importance weights, \textit{variants2.py} computes the importance weighted variants, \textit{spectrum.py} and \textit{spectrum\_hi.py} the information per mode, \textit{advection\_ratio.py} the ratio of the off diagonal to the diagonal of the tangent operator, \textit{crossover.py} the two further regimes of Section~\ref{m:sec:expns} (it takes the word \texttt{amplitude} or \texttt{posterior} as its argument), and \textit{relin.py} and \textit{relin2.py} the summary linearized once and linearized again. The numbers of Remark~\ref{m:rem:nsdiag} come from \textit{radius\_iw.py} and \textit{relin.py}, and \textit{make\_fig2.py} draws the three panel figure \textit{fig\_ns\_study.pdf} from \textit{variants.json}, \textit{radius.json} and \textit{spectrum\_hi.json}, which the book does not include.
\item[Reruns of the Navier Stokes scripts.] We reran all of these scripts, one by one, in a copy of the folder. The files that \textit{radius.py}, \textit{radius\_iw.py}, \textit{relin.py}, \textit{relin2.py}, \textit{variants.py}, \textit{variants2.py}, \textit{spectrum.py}, \textit{spectrum\_hi.py}, \textit{centring.py} and \textit{correct\_centre.py} wrote agree with the recorded files to a relative difference of at most $3\times10^{-10}$ in every entry, except for four entries of \textit{correct\_centre.json} that depend on a third derivative taken by finite differences, which differ by up to $6\times10^{-8}$. The files \textit{advection.json}, \textit{mcmc\_error.json}, \textit{crossover2.json} and \textit{crossover3.json} agree with the recorded files in every digit. In our runs \textit{hetero.py}, \textit{radius.py}, \textit{radius\_iw.py} and \textit{relin.py} took between 300 and 490 seconds, the two runs of \textit{crossover.py} about 390 and 460 seconds, \textit{relin2.py} ten minutes, and each of the others under a minute, so that the sequence took about an hour.
\item[Corrections in the Navier Stokes experiment.] The folder \textit{code/ns\_checks} holds four scripts that we wrote for the book, to measure what the corrections of Section~\ref{m:sec:loc} do in the experiment of Section~\ref{m:sec:expns}. They are run in a copy of the folder \textit{code/anc/ns}, after \textit{full.py} and \textit{hetero.py}, and write their results into that copy. The script \textit{table\_ablation.py} repeats \textit{correct\_centre.py} with the same seeds, so that it reproduces its numbers, and scores the rescaled generator with and without the shift of the center by the cubic correction and over twenty further resamplings of its draws; it gives the row of the rescaled spread of Table~\ref{m:tab:ns}, the numbers of its caption and the prior fractions of the main run, and writes \textit{table\_ablation.json}. The script \textit{regimes.py}, with the argument \texttt{amplitude} or \texttt{posterior}, repeats \textit{crossover.py} statement by statement, applies the centering and the rescaling to the same training draws, and writes \textit{regimes\_amplitude.json} or \textit{regimes\_posterior.json} with the Monte Carlo standard error of the reference and the resampling noise of the scores. The script \textit{cross\_base.py}, with the prior constant, the noise standard deviation and a tag as its arguments, runs the reference chain of the regimes for the chain seeds $77$, $78$ and $79$; we ran it at $(1.5, 0.45)$, $(2, 0.45)$ and $(3, 1.30)$. The script \textit{skew\_check.py} compares the skewness of the reference chains with that of Gaussian chains of the same length and autocorrelation. On a cluster node with an Intel Xeon Gold 6240R processor, \textit{crossover.py} followed by \textit{regimes.py} took eight to ten minutes for each regime, each run of \textit{cross\_base.py} about nine minutes and \textit{skew\_check.py} fifteen seconds, and \textit{full.py}, \textit{hetero.py}, \textit{centring.py}, \textit{correct\_centre.py}, \textit{table\_ablation.py}, \textit{mcmc\_error.py} and \textit{radius\_iw.py} took seventeen minutes together.
\item[Reconstruction of the manifold gain.] The functions \texttt{constrained} and \texttt{unconstrained} of \textit{reconstruct\_manifold\_gain.py} reconstruct the table of the paper, which has the same $\rho$ in both priors and so differs from Table~\ref{m:tab:gain} in the unconstrained row and the ratios. They calibrate the priors to the pointwise prior standard deviations $2.2$ and $2.5$ and use other cutoffs of the modes, so that they agree with the table of the paper only approximately; they take a fraction of a second. They give $20$ and $172$ effective parameters and the standard deviations and ratios of that table with differences of at most one unit in the last digit of the standard deviations, except that the ratio at $t=0.2$ comes out as $2.74$ where the paper has $2.8$; the script itself reports this cell and gives the recorded value as $2.76$. The main function of the script reads the file \textit{paper/manifold\_gain.json}, so the script as copied stops at its first line.
\item[Scripts of the paper.] The scripts \textit{fig\_heat\_stokes.py}, \textit{fig\_heat\_gbc.py}, \textit{fig\_cost.py} and \textit{fig\_nls\_bogoliubov.py} drew the figures of the paper, and \textit{make\_experiments\_tex.py} wrote the experiments section of the paper from the recorded results. Except for \textit{fig\_nls\_bogoliubov.py}, which writes \textit{fig\_nls.pdf} into the folder from which it is run, they write into a folder \textit{paper} next to the folder \textit{anc}, which the book source does not have. We created that folder in our copy and ran the five scripts, which took between 0.2 and 6 seconds each and wrote the files \textit{fig\_heat.pdf}, \textit{fig\_tg.pdf}, \textit{fig\_heat\_gbc.pdf}, \textit{fig\_cost.pdf}, \textit{fig\_nls.pdf} and \textit{experiments\_generated.tex}. Like \textit{mf\_nls.py}, \textit{fig\_nls\_bogoliubov.py} needs NumPy 2.0 or later.
\item[Earlier experiments.] The scripts \textit{exp\_study.py}, \textit{exp\_gbc\_heat.py} and \textit{exp\_gbc\_local.py} are earlier experiments with the heat equation. Their recorded results are \textit{results/study.json}, \textit{results/gbc\_heat\_T2.5\_rep1.json} and the folders \textit{rho1}, \textit{rho1e3} and \textit{rho1e5} of \textit{results/local\_hopper}, and no script of the book reads them. A full run of \textit{exp\_study.py} with its default of two hundred replicates took 32 minutes and reproduced \textit{results/study.json} to a relative difference of at most $1.1\times10^{-12}$ in every entry except the wall clock times. Its seeds are $90000+r$ for replicate $r$, $11$ for the data of the comparison of cost, $5000+M$ for the GBC run with $M$ draws and $7000+n$ for the pCN run with $n$ iterations. The first replicate of \textit{exp\_gbc\_local.py} for each of the three prior scales (twenty modes, $8000$ observations and $4000$ training draws) took 65, 92 and 130 seconds and reproduced the recorded file to a relative difference of at most $7\times10^{-11}$ in every entry above $10^{-6}$. The recorded file of \textit{exp\_gbc\_heat.py} was written by the option that makes the run quick, at $T=2.5$, and by an earlier revision of the script that also wrote two rows labeled adjoint. The present script no longer writes them, as its comments explain, so it reproduces the first four rows of the file in every digit except the time but not the later rows, which depend on the random numbers that the removed rows used.
\end{description}

\section{Rebuilding the book}
\label{sec:code_rebuild}

The book is the file \textit{main.tex} in the folder \textit{source}, which reads the chapters from \textit{chapters} and the figures from \textit{figures}. Run \texttt{pdflatex main.tex} three times in that folder; a pass takes about fifteen seconds on our laptop with pdfTeX from TeX Live 2026. The bibliography is a hand written \texttt{thebibliography} environment in \textit{chapters/bibliography.tex}, so no bibliography program is run. The passes must follow one another with no edit in between. The first pass writes the labels and the entries of the contents, and the references are undefined while it runs. The second resolves the references, which changes the length of the book and so moves the page numbers of the contents, but it typesets the contents with the page numbers of the first pass. The third uses the page numbers of the second, and a fourth pass changes neither the auxiliary file nor the contents file, as we checked in a scratch copy.

The wrapped cover is a separate document, \textit{cover/cover.tex}, built with \texttt{pdflatex} in the folder \textit{cover}. Its front image \textit{art\_ns.png} is a vorticity field of a two dimensional Navier Stokes flow on a $256\times256$ grid with viscosity $5\times10^{-4}$, drawn with the seed $1900$ by \textit{cover/art\_ns.py}, which writes the image into the folder from which it is run and takes about two minutes. The width of the spine in \textit{cover.tex} depends on the number of pages of the book and has to be set by hand.
